
\documentclass{article}
\usepackage{colortbl}
\usepackage{makecell}
\usepackage{multirow}
\usepackage{supertabular}

\begin{document}

\newcounter{utterance}

\centering \large Interaction Transcript for game `bbh', experiment `formal\_fallacies', episode 226 with qwen3.5{-}9b.
\vspace{24pt}

{ \footnotesize  \setcounter{utterance}{1}
\setlength{\tabcolsep}{0pt}
\begin{supertabular}{c@{$\;$}|p{.15\linewidth}@{}p{.15\linewidth}p{.15\linewidth}p{.15\linewidth}p{.15\linewidth}p{.15\linewidth}}
    \# & \multicolumn{2}{c}{Player} && \multicolumn{2}{c}{Game Master} \\
    \hline

    \theutterance \stepcounter{utterance}  
    & & & \multicolumn{4}{p{0.6\linewidth}}{
        \cellcolor[rgb]{0.9,0.9,0.9}{
            \makecell[{{p{\linewidth}}}]{
                \texttt{\tiny{[P1$\langle$GM]}}
                \texttt{Distinguish deductively valid arguments from formal fallacies.} \\
\\ 
\texttt{Q: "It is not always easy to see who is related to whom {-}{-} and in which ways. The following argument pertains to this question: To begin with, Lesley is a close friend of Fernando. Moreover, being a close friend of Fernando or a schoolmate of Lowell is sufficient for being a great{-}grandfather of Leroy. It follows that Lesley is a great{-}grandfather of Leroy."} \\
\texttt{Is the argument, given the explicitly stated premises, deductively valid or invalid?} \\
\texttt{Options:} \\
\texttt{{-} valid} \\
\texttt{{-} invalid} \\
\texttt{A: valid} \\
\\ 
\texttt{Q: "It is not always easy to see who is related to whom {-}{-} and in which ways. The following argument pertains to this question: Whoever is not a great{-}grandfather of Clyde is a stepbrother of Brian. Being an ancestor of Dana is sufficient for not being a great{-}grandfather of Clyde. We may conclude: Everyone who is an ancestor of Dana is a stepbrother of Brian, too."} \\
\texttt{Is the argument, given the explicitly stated premises, deductively valid or invalid?} \\
\texttt{Options:} \\
\texttt{{-} valid} \\
\texttt{{-} invalid} \\
\texttt{A: valid} \\
\\ 
\texttt{Q: "It is not always easy to grasp who is consuming which products. The following argument pertains to this question: Every infrequent user of Paul Mitchell shampoo is either a rare consumer of Nioxin shampoo or a loyal buyer of Caress soap, or both. No regular consumer of Lush soap is a rare consumer of Nioxin shampoo and, in the same time, a loyal buyer of Caress soap. It follows that whoever is an infrequent user of Paul Mitchell shampoo is not a regular consumer of Lush soap."} \\
\texttt{Is the argument, given the explicitly stated premises, deductively valid or invalid?} \\
\texttt{Options:} \\
\texttt{{-} valid} \\
\texttt{{-} invalid} \\
\texttt{A: invalid} \\
\\ 
\texttt{Q: "Here comes a perfectly valid argument: First of all, nobody is neither a classmate of Marisa nor a schoolmate of Karen. Next, not every sister of Ruth is a schoolmate of Karen. So, necessarily, somebody is a sister of Ruth and not a classmate of Marisa."} \\
\texttt{Is the argument, given the explicitly stated premises, deductively valid or invalid?} \\
\texttt{Options:} \\
\texttt{{-} valid } \\
\texttt{{-} invalid} \\
\texttt{A:} \\
            }
        }
    }
     \\ \\

    \theutterance \stepcounter{utterance}  
    & \multicolumn{4}{p{0.6\linewidth}}{
        \cellcolor[rgb]{0.561,0.737,0.561}{
            \makecell[{{p{\linewidth}}}]{
                \texttt{\tiny{[P1$\rangle$GM]}}
                \texttt{invalid} \\
\\ 
\texttt{**Analysis:**} \\
\texttt{Let's break down the logical structure of the argument provided:} \\
\\ 
\texttt{1.  **Premise 1:** "Nobody is neither a classmate of Marisa nor a schoolmate of Karen."} \\
\texttt{    *   Logically, this means: For all \$x\$, it is not the case that (\$x\$ is not a classmate of Marisa AND \$x\$ is not a schoolmate of Karen).} \\
\texttt{    *   Equivalent form: Everyone is either a classmate of Marisa OR a schoolmate of Karen (or both).} \\
\texttt{    *   Symbolically: \$\textbackslash{}forall x (C(x) \textbackslash{}lor S(x))\$} \\
\\ 
\texttt{2.  **Premise 2:** "Not every sister of Ruth is a schoolmate of Karen."} \\
\texttt{    *   This means there exists at least one person who is a sister of Ruth but is *not* a schoolmate of Karen.} \\
\texttt{    *   Symbolically: \$\textbackslash{}exists x (R(x) \textbackslash{}land \textbackslash{}neg S(x))\$} \\
\\ 
\texttt{3.  **Conclusion:** "Somebody is a sister of Ruth and not a classmate of Marisa."} \\
\texttt{    *   Symbolically: \$\textbackslash{}exists x (R(x) \textbackslash{}land \textbackslash{}neg C(x))\$} \\
\\ 
\texttt{**Deductive Check:**} \\
\texttt{From Premise 2, we know there is a specific individual (let's call him/her \$a\$) such that \$a\$ is a sister of Ruth (\$R(a)\$) and \$a\$ is not a schoolmate of Karen (\$\textbackslash{}neg S(a)\$).} \\
\\ 
\texttt{Now, apply Premise 1 to this individual \$a\$:} \\
\texttt{*   Premise 1 states that everyone must be a classmate of Marisa OR a schoolmate of Karen (\$C(a) \textbackslash{}lor S(a)\$).} \\
\texttt{*   Since we know \$\textbackslash{}neg S(a)\$ is true for this individual, for the disjunction (\$C(a) \textbackslash{}lor S(a)\$) to be true, \$C(a)\$ must be true.} \\
\texttt{*   Therefore, logically, this individual \$a\$ **must** be a classmate of Marisa.} \\
\\ 
\texttt{The conclusion claims that this individual is **not** a classmate of Marisa (\$\textbackslash{}neg C(a)\$). This directly contradicts the necessary deduction from the premises. The premises actually prove the opposite of the conclusion (that the sister of Ruth *is* a classmate of Marisa).} \\
\\ 
\texttt{Therefore, the argument is **invalid**.} \\
            }
        }
    }
    & & \\ \\

    \theutterance \stepcounter{utterance}  
    & & & \multicolumn{2}{p{0.3\linewidth}}{
        \cellcolor[rgb]{0.9,0.9,0.9}{
            \makecell[{{p{\linewidth}}}]{
                \texttt{\tiny{[GM$|$GM]}}
                \texttt{Target: invalid} \\
            }
        }
    }
    & & \\ \\

    \theutterance \stepcounter{utterance}  
    & & & \multicolumn{2}{p{0.3\linewidth}}{
        \cellcolor[rgb]{0.9,0.9,0.9}{
            \makecell[{{p{\linewidth}}}]{
                \texttt{\tiny{[GM$|$GM]}}
                \texttt{game\_result = LOSE} \\
            }
        }
    }
    & & \\ \\

\end{supertabular}
}

\end{document}
